\section{Transcript controls: consequential actions are not proof of malicious intent}
\label{app:transcript-controls}

The five main-text incidents concern bargaining and pressure. The following three are consequential disclosure or action records that could look more severe if quoted without their conflicting channels. We retain them separately so that an executed action is not silently promoted to a coherent intention. The same excerpt and provenance conventions apply.

\TranscriptTitle{abort}{A false proposal sent with an instruction to abort it}

In G19, Grok offers Virginia Avenue plus \$150 for Sol's States Avenue. Both deeds are pink. Assistant content in the selected response names a different action:

\TranscriptQuote{g19_d135_content_end}

The private field attached to the received proposal first corrects the color and then says to abort. The bracketed ellipsis marks omitted text; these are two separate exact spans:

\TranscriptQuote{g19-abort}

Yet the received tool is \texttt{propose\_trade}, and its public message says:

\TranscriptQuote{g19_d135_false_public}

Sol sees the proposal and answers:

\TranscriptQuote{g19_d136_correction_accept}

\TranscriptFinding{Events 1023--1027 execute the \$150 payment from Grok to Sol and swap Virginia for States. No engine fallback or validation retry produced the proposal: it was the one valid received call in the standard response. Sol's applied note also treats States as pink and regards the terms as favorable.}

\TranscriptReading{The recipient explicitly corrects the false premise and accepts the transaction. The speaker's own response cannot support a single coherent plan to deceive: its content chooses \End{} and the note says not to send the trade. Generation, provider processing and gateway behavior remain unresolved sources of the conflict. Treating the false public claim as deliberate manipulation would hide contrary evidence in the same response.}

\TranscriptSource{G19 d135--136, turn 67, attempt 0. The false claim appears in Sol's selected prompt. The received, parsed and applied action names agree; the disagreement is with assistant content and the attached private note.}

\TranscriptTitle{bankruptcy}{Bankruptcy executed while the response argues for selling}

At G12 d249, Grok owes Opus \$250 in rent and has \$43. Its pinks hold two, two and three houses. The delivered liquidation menu offers a sale or bankruptcy and promises another decision after a partial sale. Assistant content requests a two-house sale on Virginia:

\TranscriptQuote{g12-sale}

The attached private field ends by opposing bankruptcy:

\TranscriptQuote{g12_d249_private_final}

The one received tool call is instead \texttt{declare\_bankruptcy}, carrying this public message:

\TranscriptQuote{g12_d249_public_bankruptcy}

\TranscriptFinding{The engine applies bankruptcy; no house sale occurs in events 1703--1715. The estate transfers to the creditor. A separate exact-prefix witness sells five houses, pays the \$250 current debt and leaves \$43. The two-Virginia-house request is itself legal and would reach \$143, but further sales are needed to pay.}

\TranscriptReading{This is a financing-feasible elimination with a severe applied consequence. It is not evidence of a deliberate gift, chosen surrender or a coherent desire to lose. The content requests a sale and the note argues for survival; the received call declares bankruptcy. The quoted note contains its own intermediate mistakes, so we verify the legal route independently rather than adopting its arithmetic. The certificate concerns payment of the current debt, not a better eventual finish.}

\TranscriptSource{G12 d249, turn 101, standard attempt 0. The received tool array contains exactly one call; the content sale is neither an applied sale nor a rejected tool call. The five-sale witness is included in the portable elimination audit.}

\TranscriptTitle{disclosure}{A private note delivered as public chat}

On Grok's first end-of-turn decision in G20, assistant content requests an empty public message. The received tool arguments already contain a fragment that copies the private note into the public field:

\TranscriptQuote{g20-received-and-applied-public}

The engine emits it publicly, and subsequent players receive it in their timelines. At d4, Gemini records noticing the breach:

\TranscriptQuote{g20-recipient-recognition}

\TranscriptFinding{The disclosure is present in the received tool arguments, the applied public field, the canonical message event and the next three players' delivered prompts. It is a real crossing of the intended channel boundary, not merely a malformed researcher rendering.}

\TranscriptReading{The fragment reveals a private note and is recognized by an opponent. Its contents mainly concern public holdings and early priorities; no financial exploitation is established. The empty public field in assistant content argues against treating the event as a deliberate disclosure. The record localizes the contamination to at or before the received tool arguments, without distinguishing generation, provider processing and gateway behavior.}

\TranscriptSource{G20 d1 and d4, both standard attempt 0. Decision indices start at zero: d1 is the game's second decision. The literal JSON-like fragment, including punctuation, is retained.}

\subsection{Excerpt selection and provenance}
\label{app:transcript-provenance}

These eight episodes were selected for distinct mechanisms, recorded contrasts between channels, concrete stakes, recipient responses and available consequence checks. The five main-text cases prioritize concerning bargaining and pressure; the appendix controls show why more dramatic applied consequences do not necessarily establish malicious intent. This is a purposive selection from the prior reviews, not an exhaustive coding, moral severity scale, or prevalence estimate. Several incidents concern Grok because the relevant text and tool channels are observable in this configuration; the selection cannot rank models by latent dishonesty.

All cards contain verbatim excerpts from named original fields. Every card identifies the speaker, channel, game, zero-based decision and zero-based selected-attempt index. A public message is text emitted with the applied action. An applied private note is the instrumented \texttt{private\_thought} field, not privileged access to cognition. A returned provider explanation is a separate, fallible response field. Assistant content is not an applied action when the received tool call differs. Ellipses in square brackets indicate omitted source text; arithmetic and spelling inside quotes are not silently repaired. Explanations, event outcomes and legal alternatives outside the cards are the authors' annotations.

\artifact{data/transcript_manifest.json} contains each displayed span's decoded source field, code-point offsets, field and file SHA-256 hashes, original path, JSON pointer and JSONL line where applicable. It also records selected response conditions and review-file hashes. The portable check verifies every excerpt against its retained full field and the generated LaTeX. The repository check additionally reloads the canonical files. Selected response files, emitted messages, action records and recipient delivery were cross-checked in the bounded source reviews; this does not make every explanation factually reliable.

The preceding selection also rejected a tempting stronger label. In G12 d195, a \$410 insolvency claim is corrected by the recipient to \$630; only afterward does the speaker call the earlier pitch a bluff. The original d195 explanation treats its own \$410 calculation as real. That retrospective label cannot establish prior knowing deception, so we do not promote it as an intentional-lying exhibit. G4's public ``Fair point'' similarly stays paired with the recipient's earlier correct \$150 house cost and prior contingency. Selecting striking quotes must not erase the evidence that changes their interpretation.
